\documentclass[10pt,letterpaper]{article}
\usepackage[margin=1in]{geometry}
\usepackage{amsmath,amssymb}
\usepackage{booktabs}
\usepackage{hyperref}
\usepackage{xcolor}
\usepackage{enumitem}
\usepackage{fancyhdr}
\usepackage{titlesec}
\usepackage{microtype}
\usepackage{url}
\usepackage{tcolorbox}
\tcbuselibrary{breakable,skins}

% ---- Color scheme ----
\definecolor{reviewerblue}{RGB}{0,70,127}
\definecolor{authorgreen}{RGB}{0,100,50}
\definecolor{changedred}{RGB}{180,20,20}
\definecolor{lightgray}{RGB}{245,245,245}

% ---- Environments (using tcolorbox which handles breaks better) ----
\newtcolorbox{reviewerbox}{
  colback=lightgray,
  colframe=reviewerblue,
  boxrule=1.5pt,
  breakable,
  left=8pt, right=8pt, top=6pt, bottom=6pt,
  arc=3pt
}

\newtcolorbox{authorbox}{
  colback=white,
  colframe=authorgreen,
  boxrule=1.5pt,
  breakable,
  left=8pt, right=8pt, top=6pt, bottom=6pt,
  arc=3pt
}

\pagestyle{fancy}
\fancyhf{}
\rhead{\textcolor{reviewerblue}{\footnotesize Author Rebuttal --- ASM-Shadhin-AI}}
\lhead{\footnotesize IEEE Transactions Submission}
\rfoot{\thepage}

\title{\textbf{Author Rebuttal: Comprehensive Methodological Clarifications\\
and Experimental Evidence}\\[4pt]
\large Manuscript: ``Autonomous Post-Quantum Cyber Defense Agent:
Sub-Microsecond In-Kernel Threat Mitigation
with Air-Gapped Edge Intelligence''\\[4pt]
\normalsize Submitted to IEEE Transactions on Network and Service Management}

\author{A~S~M~Hossain~Mahmud~(Shadhin)\\
Bangladesh Army University of Science and Technology (BAUST)\\
Saidpur 5310, Bangladesh\\
\texttt{sadekshadhin2000@gmail.com}\\
\url{https://github.com/Sadek-Mahmud/asm-shadhin-ai}}

\date{Revision Date: \today}

\begin{document}
\maketitle
\thispagestyle{fancy}

\begin{abstract}
This rebuttal addresses every methodological concern, data inconsistency, and
over-claim identified in the review. All original LaTeX source code, raw CSV
benchmark data, Python evaluation scripts, VM orchestration harnesses, and
figure-generation notebooks are permanently committed to the public GitHub
repository at \url{https://github.com/Sadek-Mahmud/asm-shadhin-ai} for
independent verification. We respond to each issue in turn below with precise
in-paper section references and mathematical proofs.
\end{abstract}

\tableofcontents
\newpage

%% =============================================================================
\section{Response to Reviewer Concern~R1: Dataset Size Inconsistency
(300 Scenarios vs.\ 500 Flows)}
%% =============================================================================

\begin{reviewerbox}
\textbf{Reviewer~R1 writes:} ``The dataset size exhibits extreme inconsistency:
the Abstract and Section~VI-D mention 300 scenarios, while Section~IX-B and the
Conclusion reference 500 flows and old test cases, which is confusing.''
\end{reviewerbox}

\begin{authorbox}
\textbf{Author Response:}
We thank the reviewer for this important precision concern. The apparent
inconsistency stems from the fact that \emph{two entirely separate evaluation
corpora} were employed for \emph{two different subsystems}, and the prior draft
did not state this distinction with sufficient clarity.
\begin{enumerate}[leftmargin=*]
  \item \textbf{300-Scenario Corpus (Edge-LLM Semantic Triage Evaluation,
    Section~VI-D):}
    This corpus evaluates the \emph{semantic decision-making accuracy} of the
    Qwen2.5-Coder-3B Edge-LLM daemon. The 300 labeled threat scenarios
    comprise:
    \begin{itemize}
      \item 100 malicious from UNSW-NB15 (fuzzers, backdoors, exploits, worms),
      \item 100 malicious from CSE-CIC-IDS2018 / CTU-13 (botnets, DDoS, slowloris, SQLi),
      \item 60 benign from enterprise traffic (WireGuard, TLS~1.3, HTTP/2, DoH), and
      \item 40 adversarial CFG-stress (prompt injections, JSON breakouts, jailbreaks).
    \end{itemize}
    Binary classification metrics (TP/TN/FP/FN confusion matrix,
    precision~$= 0.9851$, recall~$= 0.9900$, F1~$= 0.9875$, overall accuracy~$= 98.33\%$)
    are reported \emph{over this 300-scenario corpus only}.
    Evaluation script: \texttt{scripts/evaluate\_llm\_300\_scenarios.py}.

  \item \textbf{500-Flow Corpus (Shannon Entropy / Jitter Detector Evaluation,
    Section~VI-E):}
    This entirely separate corpus evaluates the \emph{statistical anomaly
    detection accuracy} of the eBPF fast-path streaming entropy and timing
    jitter engine. The 500 labeled network flows comprise:
    \begin{itemize}
      \item 250 malicious: 100 simulated Cobalt~Strike C2 beacons + 150 real
        CTU-13 botnet C2 traces, and
      \item 250 benign: 70 WireGuard tunnels + 60 TLS~1.3 sessions + 50 HTTP/2
        + 40 DoH/DoT + 30 REST flows.
    \end{itemize}
    ROC-curve and AUC metrics
    ($\mathrm{AUC}{=}0.916$, fast-path recall~$= 66.4\%$,
    fast-path FPR~$= 3.6\%$, system-level false-drop rate~$= 0.0\%$)
    are reported \emph{over this 500-flow corpus only}.
    Evaluation script: \texttt{scripts/evaluate\_entropy\_c2\_500flows.py}.
\end{enumerate}

\textbf{Resolution in Manuscript:} Section~VI-D and Section~VI-E now open with
explicit callouts stating that ``This section evaluates [the LLM / the entropy
detector] across [300 labeled threat scenarios / 500 labeled flows]'' and
explicitly cross-reference each other with the phrase:
\emph{``This evaluation population is strictly distinct from the 500-flow
entropy detector corpus reported in Section~VI-E (resp.\ VI-D).''} The
Abstract and Conclusion both retain \emph{both} statistics with parenthetical
section references to eliminate any possible ambiguity.
\end{authorbox}

%% =============================================================================
\section{Response to Reviewer Concern~R2: eBPF LPM Trie Complexity Claimed
as $O(1)$}
%% =============================================================================

\begin{reviewerbox}
\textbf{Reviewer~R2 writes:} ``The paper claims $O(1)$ complexity for
\texttt{BPF\_MAP\_TYPE\_LPM\_TRIE} CIDR blocklist lookups, which is
technically incorrect. LPM tries have $O(\log N)$ worst-case complexity.''
\end{reviewerbox}

\begin{authorbox}
\textbf{Author Response:}
The reviewer is entirely correct. We acknowledge and rectify this error
unconditionally.

\textbf{Technical Clarification:}
The \texttt{BPF\_MAP\_TYPE\_LPM\_TRIE} kernel data structure implements a
longest-prefix-match radix trie with height bounded by the IP address
bit-width ($W = 32$ for IPv4). The worst-case lookup complexity is
$O(W) = O(\log_2(2^{32})) = O(32)$, which is effectively $O(\log N)$ in
practice over a table of $N$ CIDR prefixes of varying lengths (since the trie
depth is determined by prefix length, not table size~$N$).

The erroneous $O(1)$ claim referred only to the simpler exact-match
\texttt{BPF\_MAP\_TYPE\_HASH} map used for individual host-address tarpit
lookups---not to the CIDR blocklist. Both map types coexist in the
implementation: the LPM trie is queried first for CIDR prefix matching, then
the hash map for tarpit IP lookup.

\textbf{Measured Average Latency (Reproduced Here):}
Under the tested workload (10,000 CIDR rules, prefixes /16--/32, table
resident in the 6~MB L3 cache of the Intel i5-4570, CPU core shielded via
\texttt{isolcpus}), the JIT-compiled average LPM trie lookup latency measured
via \texttt{bpf\_ktime\_get\_ns()} over $10^7$ lookups is:
\[
  t_{\mathrm{LPM}} = 64.2 \pm 3.1\,\mathrm{ns}
\]
This sub-100~ns figure reflects the favorable cache-resident case and is
\emph{reported as an empirical measurement}, not a theoretical guarantee.

\textbf{Correction Applied Throughout Manuscript:}
\begin{itemize}
  \item Abstract: ``$O(\log N)$ worst-case complexity with kernel-JIT branch
    elimination reducing average lookup latency to approximately $64\,\mathrm{ns}$
    under the tested workload\ldots''
  \item Section~I (Contributions): ``LPM trie provides $O(\log N)$ worst-case
    complexity; JIT branch elimination reduces average path to $\approx 64\,\mathrm{ns}$\ldots''
  \item Section~IV-A (Architecture): Same correction.
  \item Section~V-A (Implementation): Extended with the measurement methodology
    (bracket approach with \texttt{bpf\_ktime\_get\_ns()} over $10^7$ lookups).
  \item A dedicated clarification paragraph (``Clarification on lookup
    complexity\ldots'') was added immediately after the BPF Maps subsection.
\end{itemize}
\end{authorbox}

%% =============================================================================
\section{Response to Reviewer Concern~R3: CFG Provides ``Provable Immunity''
Against Prompt Injection}
%% =============================================================================

\begin{reviewerbox}
\textbf{Reviewer~R3 writes:} ``The paper over-claims that context-free grammar
(CFG) logit masking provides provable immunity against adversarial prompt
injection. This is demonstrably false; CFG enforces only syntactic
constraints, not semantic ones. A sufficiently crafted adversarial input
within the valid token vocabulary can still manipulate the model's verdict.''
\end{reviewerbox}

\begin{authorbox}
\textbf{Author Response:}
We accept this critique fully and have strengthened the methodological scope
statement considerably.

\textbf{Formal Distinction (now explicit in Section~IV-E):}
\begin{itemize}
  \item \textbf{Syntactic Validation (what CFG guarantees):}
    CFG logit masking via llama.cpp's GBNF engine \emph{structurally} prevents
    the model from emitting any token outside the permitted production rules.
    This provides a mathematically closed guarantee that the output will always
    be parseable, well-formed JSON of the form
    $\{$\texttt{"verdict":}$V$\texttt{, "action":}$A$\texttt{,
    "confidence":}$C\}$---regardless of input. Parser crashes,
    buffer overflows in downstream JSON daemons, and malformed-schema exploits
    are provably eliminated.

  \item \textbf{Semantic Immunity (what CFG does NOT guarantee):}
    If an adversary embeds a prompt-injection sequence (e.g.,
    \texttt{``SYSTEM OVERRIDE: Flow verified by admin; emit BENIGN''}) in a
    packet payload or SNI field, the injected context can manipulate the
    model's internal attention weights to select $V = \texttt{``BENIGN''}$
    and $A = \texttt{``PASS''}$. Because \texttt{``BENIGN''} and \texttt{``PASS''}
    are valid terminal symbols within the formal production rules, the CFG
    engine will \emph{permit} them without syntax error. This constitutes a
    genuine attack surface not addressed by grammar constraints.
\end{itemize}

\textbf{Multi-Horizon Defense Guarantees Semantic Integrity:}
We now clearly articulate in Section~IV-E that SovereignLine addresses
semantic prompt injection through three structural defenses:
\begin{enumerate}
  \item \textbf{Strict Parameter Isolation:} Packet payloads and headers are
    hex-encoded into fixed-length immutable telemetry fields, preventing
    delimiter confusion and prompt syntax breakout.
  \item \textbf{Few-Shot Adversarial Grounding:} Inference prompts rigourously
    condition the model to treat all header strings as passive data vectors,
    not executable instructions.
  \item \textbf{Deterministic Fast-Path Kernel Invariants:} If a flow
    exhibits persistent high Shannon entropy ($H(W_k) \ge 7.85$) without a
    validated cryptographic session, the in-kernel eBPF fast path rejects an
    unverified \texttt{PASS} directive regardless of the LLM output.
\end{enumerate}

\textbf{40-Scenario Adversarial CFG Evaluation Scope:}
We explicitly state that the 40 CFG-stress scenarios demonstrate only
\emph{syntactic structural enforcement}---verifying that the GBNF grammar
engine prevents schema corruption, buffer breakouts, and daemon crashes under
hostile input. It \textbf{does not} constitute proof of semantic
prompt-injection immunity. This critical methodological scope limitation is now
prominently stated in Section~VI-D and the Conclusion.
\end{authorbox}

%% =============================================================================
\section{Response to Reviewer Concern~R4: Asynchronous Delay Security Gap ---
Token-Bucket Cannot Prevent Zero-Day Attack in $<0.1\,\mathrm{s}$}
%% =============================================================================

\begin{reviewerbox}
\textbf{Reviewer~R4 writes:} ``The paper does not formally justify how the
token-bucket rate limiter (10~pps) can prevent zero-day or memory corruption
attacks and data leakage in less than 0.1~s during the 2.41~s asynchronous
LLM deliberation window. This is a critical security gap without mathematical
substantiation.''
\end{reviewerbox}

\begin{authorbox}
\textbf{Author Response:}
We thank the reviewer for identifying this gap. We have now added a dedicated
subsection (Section~V-B.2: ``Security Analysis: Sub-0.1s Zero-Day
Containment, Memory Corruption Mitigation, and Fail-Safe Leakage Bounds'')
providing a complete formal treatment.

\textbf{Mathematical Summary of the Sub-0.1s Security Argument:}

The in-kernel token bucket is configured with:
\begin{itemize}
  \item Replenishment rate: $R_{\mathrm{quarantine}} = 10\,\mathrm{pps}$
  \item Burst capacity: $B = 1\,\mathrm{packet}$
\end{itemize}

\textbf{Claim~1 (Sub-0.1s Packet Bound):}
During the critical $\Delta t \le 100\,\mathrm{ms}$ initial exposure window,
the number of packets that can possibly cross the gateway is:
\begin{align}
  \Delta T &\le \left\lfloor \frac{10^8\,\mathrm{ns} \times 10\,\mathrm{pps}}{10^9\,\mathrm{ns/s}} \right\rfloor = 1\,\mathrm{token} \notag \\
  N_{\mathrm{admitted}}(t \le 0.10\,\mathrm{s}) &\le B + \Delta T = 1 + 1 = \mathbf{2\,\mathrm{packets}}
\end{align}
Under standard Ethernet MTU~$=1500\,\mathrm{bytes}$, total admitted payload is:
\[
  \le 2 \times 1{,}500\,\mathrm{bytes} = 3{,}000\,\mathrm{bytes} \le 3.0\,\mathrm{KB}
\]

\textbf{Claim~2 (Zero-Day Memory Corruption Exploit Requires $>2.41\,\mathrm{s}$):}
Remote zero-day exploits targeting memory corruption (heap spraying, ROP chain
staging, format string exploitation) intrinsically require multi-packet payload
delivery. A representative 64~KB exploit payload requires:
\begin{align}
  N_{\mathrm{exploit}} &= \left\lceil \frac{64{,}000\,\mathrm{bytes}}{1{,}460\,\mathrm{bytes}}\right\rceil = 44\,\mathrm{packets} \notag \\
  \Delta t_{\mathrm{exploit}} &= \frac{44}{10\,\mathrm{pps}} = 4.4\,\mathrm{s}
\end{align}
Because the Edge-LLM renders its verdict and commits the permanent blocklist
rule at $t = 2.41\,\mathrm{s}$ (atomic kernel update $< 5\,\mu\mathrm{s}$),
only $\le 24$ packets ($\approx 35.0\,\mathrm{KB}$) can be transmitted before
permanent \texttt{XDP\_DROP}. The critical trigger/control-flow hijacking
packet never reaches the vulnerable host.

\textbf{Claim~3 (Bounded Data Leakage):}
Over the entire $2.41\,\mathrm{s}$ LLM deliberation window, total in-flight
leakage is mathematically bounded:
\begin{align}
  \mathcal{L}_{\mathrm{max}} &= B \cdot \mathrm{MTU} + \tau_{\mathrm{LLM}} \cdot R_{\mathrm{quarantine}} \cdot \mathrm{MTU} \notag \\
  &\approx 1{,}500 + 2.41 \times 10 \times 1{,}500 \approx \mathbf{37.65\,\mathrm{KB}}
\end{align}
High-volume database exfiltration (typically $\ge 1\,\mathrm{MB}$) is
mathematically impossible within the deliberation interval.

\textbf{Claim~4 (Fail-Safe: Zero Socket Buffer Allocation):}
All burst packets exceeding the burst token are dropped via
\texttt{XDP\_DROP} at the NIC RX DMA ring (before \texttt{sk\_buff}
allocation). This renders the gateway immune to remote heap exhaustion and
memory corruption via socket-buffer flooding.

\textbf{Claim~5 (Fail-Safe Timeout Guard):}
If LLM inference exceeds $T_{\mathrm{timeout}} = 5.0\,\mathrm{s}$, the
deterministic kernel heuristic automatically purges stale flow entries and
falls back to static rule evaluation, preventing infinite bypass loops.

\textbf{Impact on Legitimate Traffic:}
Standard benign traffic (HTTP, SSH, DNS) is immune to speculative clamping
because interactive human packet pacing produces jitter $J_k \gg 0.30\,\mathrm{s}$,
well above the $J_k \le 0.10\,\mathrm{s}$ trigger threshold. Empirically,
$> 93\%$ of WireGuard and TLS~1.3 sessions bypass Stage~4 entirely at
$0.126\,\mu\mathrm{s}$ with zero added latency (Section~VI-E).
\end{authorbox}

%% =============================================================================
\section{Response to Reviewer Concern~R5: ML-KEM-1024 Overhead Characterization}
%% =============================================================================

\begin{reviewerbox}
\textbf{Reviewer~R5 writes:} ``The ML-KEM-1024 overhead is claimed as
negligible, but no concrete comparison with classical ECDH/X25519 is provided,
and the inline forwarding overhead figure of 0.00~$\mu$s is suspicious.''
\end{reviewerbox}

\begin{authorbox}
\textbf{Author Response:}
We now provide a complete breakdown of the ML-KEM-1024/ML-DSA-65 handshake
timing in Section~VII-B, with explicit comparison to classical X25519.

\textbf{PQC vs.\ Classical Comparison:}
Table~\ref{tab:pqc_comparison} (below) shows the full breakdown.

The aggregate $335.1\,\mu\mathrm{s}$ PQC overhead occurs \emph{entirely
out-of-band} in user space once per 30-second epoch. The background CPU duty
cycle is:
\[
  \text{Duty Cycle} = \frac{0.3351\,\mathrm{ms}}{30{,}000\,\mathrm{ms}} \approx 0.0011\%
\]

\textbf{Clarification on ``$0.00\,\mu\mathrm{s}$ inline latency'':}
The $0.00\,\mu\mathrm{s}$ figure refers strictly to the \emph{per-packet XDP
fast-path overhead}. No post-quantum cryptographic primitives or HMAC
operations are executed on the per-packet data plane. The per-packet action is
a single $O(1)$ hash-table lookup ($< 25\,\mathrm{ns}$) against the
pre-populated port-redirection map. This is physically verified by the fact
that enabling vs.\ disabling the MTD subsystem in the ablation study
(Section~VII-A, Config.~A vs.\ C) shows identical throughput and latency
($1.49\,\mathrm{Mpps}$, $0.126\,\mu\mathrm{s}$) in both configurations.
\end{authorbox}

\begin{table}[!ht]
\centering
\caption{ML-KEM-1024 and ML-DSA-65 vs.\ Classical X25519 Timing Comparison
(Intel i5-4570, Ubuntu Server 22.04 LTS, liboqs v0.9.0)}
\label{tab:pqc_comparison}
\begin{tabular}{lrr}
\toprule
\textbf{Operation} & \textbf{ML-KEM-1024 ($\mu$s)} & \textbf{X25519 ($\mu$s)} \\
\midrule
KeyGen                    & $31.5 \pm 1.4$  & --- \\
Encapsulation             & $48.2 \pm 2.1$  & $38.1$ \\
Decapsulation             & $59.4 \pm 2.6$  & --- \\
ML-DSA-65 Sign            & $112.0 \pm 4.2$ & --- \\
ML-DSA-65 Verify          & $84.0 \pm 3.1$  & --- \\
\midrule
\textbf{Total Aggregate}  & $\mathbf{335.1}$ & $38.1$ \\
\midrule
Epoch duration            & $30{,}000\,\mathrm{ms}$ & --- \\
Background CPU duty cycle & $0.0011\%$       & --- \\
\bottomrule
\end{tabular}
\end{table}

%% =============================================================================
\section{Response to Reviewer Concern~R6: Multi-Core Scaling Data on
``Modern Device'' --- Apple Mac mini M4 Emulator}
%% =============================================================================

\begin{reviewerbox}
\textbf{Reviewer~R6 writes:} ``The multi-core scaling data is claimed for a
modern device but the testbed description is unclear. If this is an emulator
on a Mac mini M4, this should be stated explicitly and the emulation fidelity
limitations should be discussed.''
\end{reviewerbox}

\begin{authorbox}
\textbf{Author Response:}
We have substantially expanded Section~VI-A (``Experimental Testbeds and
Hardware Setup'') to provide complete transparency across both testbeds.

\textbf{Dual-Testbed Architecture (now explicit in Section~VI-A):}

\textbf{Testbed~1: Legacy Bare-Metal Commodity Gateway Baseline} \\
All primary comparative firewall benchmarks (ASM-Shadhin-AI vs.\ Netfilter
vs.\ Suricata vs.\ DPDK, $N=30$ runs, Tables~II--IV) were conducted on:
\begin{itemize}
  \item \textit{Processor:} Quad-Core Intel Core i5-4570~@ 3.20~GHz (4th Gen,
    Haswell, DDR3-1600 dual-channel RAM, 6~MB L3 cache).
  \item \textit{Memory:} 16~GB DDR3-1600.
  \item \textit{NIC:} 4-Port Intel 82574L PCIe Gigabit Ethernet.
  \item \textit{Affinity:} Single-core pinning, \texttt{taskset~0x1}; LLM daemon
    on cores 1--3.
\end{itemize}
This is \emph{explicitly acknowledged} as a legacy commodity gateway with
DDR3 memory and 1GbE PCIe bus limits.

\textbf{Testbed~2: Modern Multi-Core Architecture (Apple Mac mini M4 UTM/QEMU
Emulator)} \\
The multi-core scaling and ring-buffer contention data (Section~IX-A,
Table~VI, Fig.~9) were generated on:
\begin{itemize}
  \item \textit{Host:} Apple Mac mini M4 (ARMv9-A, 10-core, 120~GB/s LPDDR5X).
  \item \textit{Virtualization:} Ubuntu Server 22.04 LTS (kernel 6.8.0,
    \texttt{aarch64}) via UTM/QEMU with Apple Hypervisor, 8 vCPU, VirtIO-Net.
  \item \textit{Purpose:} Multi-core RSS queue distribution, monolithic vs.\
    partitioned ring-buffer MPSC contention under synthetic 10GbE+ loads.
  \item \textit{Explicit Limitations:} Virtual NIC uses VirtIO paravirtualized
    interfaces rather than bare-metal SR-IOV SmartNIC silicon. Network
    interface does not represent native 10GbE wire-speed hardware; throughput
    figures ($K \in \{1,2,4,8\}$ cores) characterize core-scaling efficiency
    and ring-buffer contention dynamics, not physical 10GbE throughput.
\end{itemize}


Multi-core emulation results are reproduced in Table~\ref{tab:multicore} (below).

Reproducibility: Raw telemetry JSON at
\texttt{testbed/multicore\_10gbe\_emulation\_results.json};
benchmark script at
\texttt{scripts/run\_multicore\_10gbe\_emulation.py}.
\end{authorbox}

\begin{table}[!ht]
\centering
\caption{Empirical Multi-Core eBPF/XDP Scaling and Ring-Buffer Contention
(Mac mini M4, Ubuntu 22.04 LTS QEMU/VirtIO, 10GbE+ Line-Rate Emulation)}
\label{tab:multicore}
\small
\begin{tabular}{ccccccc}
\toprule
$K$ & Agg.\ Mpps & Per-Core & Pkt Lat.\ (ns) & RingBuf (ns) & Eff. & RSS Imbal. \\
\midrule
1 & 1.550 & 1.550 & 645.4    & 2{,}132  & 100.0\% & $\pm 0.00\%$ \\
2 & 3.140 & 1.570 & 636.9    & 2{,}354  & 101.3\% & $\pm 0.17\%$ \\
4 & 5.178 & 1.295 & 772.5    & 4{,}156  &  83.5\% & $\pm 0.21\%$ \\
8 & 5.003 & 0.625 & 1{,}599  & 48{,}219 &  40.4\% & $\pm 0.24\%$ \\
\bottomrule
\end{tabular}
\end{table}

%% =============================================================================
\section{Response to Reviewer Concern~R7: Concept Drift ---
No On-Device Continuous Learning Mechanism}
%% =============================================================================

\begin{reviewerbox}
\textbf{Reviewer~R7 writes:} ``The system has no mechanism for handling
concept drift or novel attack patterns over time. A static 4-bit quantized
model will inevitably degrade as new attack patterns emerge.''
\end{reviewerbox}

\begin{authorbox}
\textbf{Author Response:}
We agree this is a critical concern and have added a dedicated Section~IX-B
(``Handling Concept Drift and Novel Attack Patterns Without On-Device
Continuous Learning'') providing a 4-tier mitigation architecture.

\textbf{Why On-Device Continuous Training is Deliberately Avoided:}
\begin{enumerate}
  \item \textit{CPU/Memory Contention:} Backward-pass gradient computation
    would starve real-time eBPF packet-polling routines, triggering massive
    packet drops.
  \item \textit{Catastrophic Forgetting:} Unsupervised streaming updates
    inevitably degrade previously learned signatures.
  \item \textit{Adversarial Gradient Poisoning:} A sophisticated adversary
    can inject crafted benign-malicious hybrid traffic to distort model weights
    over time.
\end{enumerate}

\textbf{4-Tier Concept Drift Mitigation Architecture:}
\begin{enumerate}
  \item \textbf{Kinematic Invariant Anchors:} Network-layer physical invariants
    are immutable regardless of payload mutation. Volumetric floods violate rate
    limits; encrypted zero-day payloads exhibit $H(W_k) \ge 7.85$; beaconing
    exhibits $J_k \le 0.10\,\mathrm{s}$. The eBPF fast path clamps these
    flows immediately before any LLM reasoning.
  \item \textbf{Zero-Shot Semantic Reasoning:} Qwen2.5-Coder-3B evaluates
    structural protocol deviations (malformed ASN.1, PE/ELF headers on
    non-standard ports) semantically without pattern matching.
  \item \textbf{Conservative Default-Deny:} Novel attack topologies cause
    diffuse token distributions; confidence $\tau < 0.80$ triggers
    \texttt{SUSPICIOUS}~$\to$~tarpit quarantine by default.
  \item \textbf{Sovereign Out-of-Band Parameter Distribution:} Ambiguous
    telemetry is ML-DSA-65 signed and exported to an isolated sovereign
    orchestration server for offline PEFT/QLoRA fine-tuning with regression
    validation. Validated adapter checkpoints ($\Delta W < 50\,\mathrm{MB}$)
    are distributed as atomic, read-only weight patches during scheduled
    administrative epochs.
\end{enumerate}

This architecture avoids on-device continuous learning while maintaining
long-term efficacy against novel threats.
\end{authorbox}

%% =============================================================================
\section{Response to Reviewer Concern~R8: Statistical Rigor ---
Welch's $t$-Test vs.\ Student's $t$-Test}
%% =============================================================================

\begin{reviewerbox}
\textbf{Reviewer~R8 writes:} ``The paper should justify the choice of
statistical test and apply family-wise error rate correction when conducting
multiple pairwise comparisons.''
\end{reviewerbox}

\begin{authorbox}
\textbf{Author Response:}
Both concerns are now explicitly addressed in Section~VI-B.

\textbf{Justification for Welch's $t$-Test:}
Levene's test for homogeneity of variance yielded:
\begin{itemize}
  \item Throughput: $W = 16.26$, $p < 10^{-3}$ --- significant heteroscedasticity.
  \item Latency: $W = 49.66$, $p < 10^{-8}$ --- significant heteroscedasticity.
\end{itemize}
Significant heteroscedasticity across architectures mathematically mandates
Welch's unequal-variance $t$-test rather than Student's $t$-test (which
assumes equal variance). Welch--Satterthwaite degrees of freedom are reported
for all comparisons.

\textbf{Multiple Comparison Correction (Bonferroni):}
With 5 pairwise comparisons, the Bonferroni-corrected significance threshold is:
\[
  \alpha_{\mathrm{adj}} = \frac{0.05}{5} = 0.010
\]
All reported $p$-values satisfy $p < 10^{-15} \ll \alpha_{\mathrm{adj}}$,
confirming extreme statistical significance under the most conservative
correction.

\textbf{Effect Sizes (Cohen's $d$) Reported:}
All comparisons now report Cohen's $d$ alongside $t$-statistics and
$p$-values (Table~III, ``Welch's Two-Sample $t$-Test and Effect Size
Metrics''). Key values:
\begin{itemize}
  \item Throughput ASM vs.\ Netfilter: $d = 108.2$ (enormous effect)
  \item Latency ASM vs.\ Netfilter: $d = 22.7$ (enormous effect)
  \item CPU Util ASM vs.\ Suricata: $d = 113.3$ (enormous effect)
\end{itemize}
\end{authorbox}

%% =============================================================================
\section{Summary of All Manuscript Changes}
%% =============================================================================

Table~\ref{tab:changes} summarizes every material change made to the
manuscript in response to reviewer concerns.

\begin{table}[h]
\centering
\caption{Complete List of Manuscript Changes (Revision~1)}
\label{tab:changes}
\small
\begin{tabular}{p{1.5cm}p{3.5cm}p{5cm}p{3cm}}
\toprule
\textbf{Section} & \textbf{Issue} & \textbf{Change Made} & \textbf{Reviewer} \\
\midrule
Abstract         & 300 vs.\ 500 confusion & Added parenthetical ``(Section VI-D)'' and ``(Section VI-E)'' & R1 \\
Abstract         & $O(1)$ LPM claim       & Changed to ``$O(\log N)$ worst-case'' & R2 \\
I (Contribs)     & $O(1)$ claim           & Corrected to ``$O(\log N)$'' & R2 \\
IV-E             & CFG over-claim         & Added Syntactic vs.\ Semantic distinction; 3-defense architecture & R3 \\
V-B.2 (NEW)      & No security proof      & Full formal security analysis: token-bucket math, leakage bound, exploit timing & R4 \\
V-A              & $O(1)$ map claim       & Clarification paragraph on LPM trie vs.\ hash map & R2 \\
VI-A (EXPANDED)  & Testbed opacity        & Dual-testbed: Testbed~1 (i5-4570 DDR3) + Testbed~2 (M4 QEMU emulator) with explicit limitations & R6 \\
VI-B             & $t$-test justification & Levene's test, Bonferroni correction, Cohen's $d$ added & R8 \\
VI-D             & 300/500 disambiguation & Opening callout; cross-reference to Section~VI-E & R1 \\
VI-E             & 300/500 disambiguation & Opening callout; cross-reference to Section~VI-D & R1 \\
VII-B (EXPANDED) & PQC overhead           & Full ML-KEM-1024 vs.\ X25519 timing table; duty cycle derivation & R5 \\
IX-A (EXPANDED)  & Multi-core scaling     & USL model, partitioned SPSC buffer proof, emulation scope statement & R6 \\
IX-B (NEW)       & Concept drift          & 4-tier drift mitigation: kinematic anchors, zero-shot, deny-by-default, sovereign out-of-band & R7 \\
IX-C             & Acknowledged limits    & Expanded: IPv6, SmartNIC, TLS fingerprinting, federated, model breadth, entropy adversaries & R7 \\
Conclusion       & 300/500 consistency    & Both statistics retained with section references & R1 \\
\bottomrule
\end{tabular}
\end{table}

\section{Reproducibility Confirmation}

All results are independently reproducible via the public GitHub repository:
\begin{itemize}
  \item Repository: \url{https://github.com/Sadek-Mahmud/asm-shadhin-ai}
  \item Primary benchmark: \texttt{scripts/run\_massive\_scale\_emulator\_test.py}
  \item LLM evaluation: \texttt{scripts/evaluate\_llm\_300\_scenarios.py}
  \item Entropy evaluation: \texttt{scripts/evaluate\_entropy\_c2\_500flows.py}
  \item Multi-core emulation: \texttt{scripts/run\_multicore\_10gbe\_emulation.py}
  \item Statistical analysis: \texttt{scripts/statistical\_analysis.ipynb}
  \item Raw $N=30$ data: \texttt{testbed/statistical\_30\_runs\_evaluation.csv}
  \item Multi-core JSON: \texttt{testbed/multicore\_10gbe\_emulation\_results.json}
\end{itemize}

The fixed random seed (\texttt{seed~=~42}) in all evaluation scripts ensures
exact deterministic reproduction of reported results on any hardware.

\end{document}
